一个\emph{集合}是一些对象的汇集，这些对象称为该集合的\emph{元素}。
若 $x$ 是 $A$ 的一个元素，则记作 $x \in A$。
集合是\emph{外延的}——它们完全由其元素决定。
集合可以通过明确\emph{列举}其元素来确定，也可以通过给出其元素共同具有的某个性质来确定（\emph{抽象}）。
外延性意味着，列举或确定集合元素时的顺序或方式无关紧要。
要证明 $A$ 和 $B$ 是同一个集合（即 $A = B$），就必须证明 $A$ 的每个元素都是 $B$ 的元素，反之亦然。

一些重要的集合包括自然数（$\Nat$）、整数（$\Int$）、有理数（$\Rat$）和实数（$\Real$），也包括对象的\emph{字符串}（$X^*$）和无限\emph{序列}（$X^\omega$）。
如果 $A$ 的每个元素同时也是 $B$ 的元素，那么 $A$ 就是 $B$ 的\emph{子集}，记作 $A \subseteq B$。
一个集合 $B$ 的所有子集的全体本身也构成一个集合，称为 $B$ 的\emph{幂集} $\Pow{B}$。
我们可以构造集合的\emph{并集} $A \cup B$ 与\emph{交集} $A \cap B$。
一个\emph{有序对} $\tuple{x, y}$ 由两个对象 $x$ 和 $y$ 组成，但具有特定的顺序。
有序对 $\tuple{x, y}$ 和 $\tuple{y, x}$ 是不同的对（除非 $x = y$）。
所有满足 $x \in A$ 且 $y \in B$ 的有序对 $\tuple{x, y}$ 构成的集合，称为 $A$ 与 $B$ 的\emph{笛卡尔积} $A \times B$。
我们用 $A^2$ 表示 $A \times A$；例如，$\Nat^2$ 就是全体自然数对构成的集合。